\section{The multiplicative group of complex units}
\label{sec:units}

\subsection{Algebraic units as a subtype}
For a commutative ring $A$, Cubical's $A^{\times}$ is the predicate of
elements equipped with multiplicative-inverse evidence.  The local definition
is the dependent sum
\begin{equation}
 \CHunits\coloneqq\sum_{z:\CH}(z\in\CH^{\times}).
 \label{eq:units-sigma}
\end{equation}
An element is a complex number together with algebraic evidence that it is a
unit.  This uses neither a quotient nor a classical coordinate test for
nonzero values.
\formal{\texttt{src/ConstructiveComplex/Units.agda}}
{\texttt{ComplexUnits}}
{The definition is \texttt{Σ Complex U.Rˣ}, where \texttt{U = Units ComplexCommRing}.}

Membership in the unit predicate is propositional by \agda{∈-isProp}.
Therefore two unit records with equal first projections are equal, by
\agda{Σ≡Prop}; and since $\CH$ is a set, the sigma type is a set by
\agda{isSetΣSndProp}.
\upstream{\texttt{Cubical/Foundations/Powerset.agda}}{\texttt{∈-isProp}}
\upstream{\texttt{Cubical/Foundations/HLevels.agda}}{\texttt{isSetΣSndProp}; \texttt{Σ≡Prop}}
\formal{\texttt{src/ConstructiveComplex/Units.agda}}
{\texttt{complexUnitsIsSet}; \texttt{unitPath}}
{The private \texttt{unitPath} lifts equality of underlying complex numbers to equality of units.}

\subsection{Group operations and proof}
Writing a unit as $(z,p)$, the source defines
\begin{align}
1_{\CHunits}&=(1,\text{unit evidence}),\\
(z,p)(w,q)&=(zw,\text{closure evidence}),\\
(z,p)^{-1}&=(z^{-1},\text{closure evidence}).
\end{align}
The generic \agda{Units} API supplies the evidence
\agda{RˣContainsOne}, \agda{RˣMultClosed}, and \agda{RˣInvClosed}; its
\agda{·-rinv} establishes $zz^{-1}=1$ for a certified unit.
\upstream{\texttt{Cubical/Algebra/CommRing/Properties.agda}}{\texttt{Units.RˣContainsOne}; \texttt{Units.RˣMultClosed}; \texttt{Units.RˣInvClosed}; \texttt{Units.·-rinv}}
\formal{\texttt{src/ConstructiveComplex/Units.agda}}
{\texttt{complexUnitOne}; \texttt{complexUnitMul}; \texttt{complexUnitInv}}
{These definitions instantiate the generic unit API at \texttt{ComplexCommRing}.}

\begin{proposition}
$\CHunits$ is a group under the displayed multiplication.
\end{proposition}
\begin{proof}
For $u=(z,p)$, $v=(w,q)$, and $r=(t,s)$, associativity becomes
$z(wt)=(zw)t$ after \agda{Σ≡Prop}, hence is
\agda{complexMulAssociative}.  The right unit law similarly becomes
\agda{complexMulUnitRight}, and the generic inverse theorem gives the right
inverse.  Cubical's \agda{makeGroup-right} packages sethood and these three
right-handed laws into a group; its record supplies the derived left unit and
left inverse laws.
\end{proof}
\formal{\texttt{src/ConstructiveComplex/Units.agda}}
{\texttt{complexUnitMulAssociative}; \texttt{complexUnitMulUnitRight}; \texttt{complexUnitMulInverseRight}; \texttt{ComplexUnitsGroup}; \texttt{complexUnitMulUnitLeft}; \texttt{complexUnitMulInverseLeft}}
{The first three are inputs to \texttt{makeGroup-right}; the latter two are derived from the resulting \texttt{Group} record.}

\begin{warning}[Boundary of the unit construction]
No declaration proves an equivalence between algebraic unit membership and a
coordinate-level apartness or nonzero condition on $(a,b)$.  In particular,
the formalization does not claim a constructive coordinate criterion for
membership in $\CHunits$.
\end{warning}
